In this section, we will prove Theorem \ref{thm:intro-main} assuming two results from later in the paper.
Recall our standing assumption that $n \ge 3$ is a positive integer.  In Sections \ref{sec:ThomPushout}-\ref{sec:Toda}, we studied the unit map
$$\pi_{8n-1} \bS \to \pi_{8n-1} \MO \langle 4n \rangle.$$
In Lemma \ref{lem:unit-surjectivity}, we showed that this map is surjective.  In Lemma \ref{lem:Reduction1}, we showed that the subgroup $\mathcal{J}_{8n-1}$ is in the kernel of this map.\footnote{We remark that this and the preceeding statement are not original to us and may be proven using classical tools of geometric topology---see \Cref{rmk:Qdef}} Furthermore, modulo $\mathcal{J}_{8n-1}$, every element in the kernel is an integer multiple of a single class, which by Lemma \ref{lem:toda-main} is given by the Toda bracket $w$.

Our task here is to show that, for $n \geq 32$, this element $w$ is trivial modulo $\mathcal{J}_{8n-1}$. Our strategy will be to prove, separately for each prime number $p$, that $w$ is trivial after $p$-localization. 

\begin{thm} \label{thm:w=0}
  Fix a prime number $p$.  The element
  $$w \in (\pi_{8n-1} \bS) / \J_{8n-1}$$
  is $p$-locally trivial if any of the following conditions are met:
  \begin{itemize}
  \item $p > 3$.
  \item $n \geq 32$ and $p=3$.
  \item $n \geq 17$ and $p=2$.
  \end{itemize}
\end{thm}

\begin{cnv}
For the rest of this section we will work $p$-locally for a fixed prime number $p$.  For example, we use $\pi_*\Ss$ to denote $\pi_*\Ss_{(p)}$.
\end{cnv}

The proof that $w \in \mathcal{J}_{8n-1}$ will proceed by using two results from later in the paper.   These results respectively
\begin{enumerate}
    \item establish a lower bound on the $\HFp$-Adams filtration of $w$, and
    \item exhibit an upper bound on the $\HFp$-Adams filtrations of elements of $\coker (J)$.
\end{enumerate}
To explain further, we recall the following definition, which appeared in the statement of \Cref{thm:Burklund}:

\begin{dfn} \label{dfn:h}
  For each prime number $p>2$ and each integer $k>0$, let $\Gamma_p(k)$ denote the minimal $m$ such that every $\alpha \in \pi_{k}\Ss_{(p)}$ with $\HFp$-Adams filtration strictly greater than $m$ is in the image of $J$.
  Similarly, let $\Gamma_2(k)$ denote the minimal $m$ such that every $\alpha \in \pi_{k}\Ss_{(2)}$ with $\HFt$-Adams filtration strictly greater than $m$ is in the subgroup generated by the image of $J$ and the $\mu$-family.
\end{dfn}

\begin{rmk}
At $p=2$, the elements of $\pi_*\Ss$ in the $\mu$ family are of degrees $1$ and $2$ modulo $8$, and in particular do not occur in degree $8n-1$.  Thus, an element in $\pi_{8n-1}\Ss_{(p)}$ with $\HFt$-Adams filtration greater than $\Gamma_p(8n-1)$ is automatically in the image of $J$.
\end{rmk}

If we let $AF(w)$ denote the $\HFp$-Adams filtration of (some choice of) $w$, then it will suffice to show that
\begin{align}\label{eqn:BIG} \Gamma_p (8n-1) < AF(w). \end{align}
\Cref{sec:SyntheticToda} will be devoted to establishing a lower bound on $AF(w)$. To state this bound, we establish additional notation.

\begin{dfn}
We define the integer $N_2$ by the formula
\[N_2 = h(4n-1) - \lfloor \log_2 (8n) \rfloor + 1,\]
where $h(k)$ is the number of integers $0 < s \leq k$ which are congruent to $0,1,2$ or $4$ mod $8$.
For an odd prime $p$, we define 
\[N_p = \left\lfloor \frac{4n}{2p-2} \right\rfloor - \left\lfloor \log_p (4n) \right\rfloor.\]
\end{dfn}

\begin{thm}[Proven as \Cref{thm:AdamsBound}]\label{thm:filt-bound-finale}
There exists a choice of $f$ in the statement of Lemma \ref{lem:toda-main} such that the $\mathrm{H}\F_p$-Adams filtration of the Toda bracket $w$ is at least $2N_p-1$.
\end{thm}

%% \begin{proof}
%% This will be proven as \Cref{thm:AdamsBound}.
%% \end{proof}

\begin{rmk}
  Our use of Adams filtrations in the proof of \Cref{thm:w=0} is inspired by arguments of Stolz in \cite{StolzBook}. 
  In partiular, it follows from a theorem of Stolz that there is a lower bound of size $N_2$ on the $\HFt$-Adams filtration of $w$ \cite[Satz 12.7]{StolzBook}.
  The doubling of Stolz's lower bound is one of the main contributions of this paper.
\end{rmk}

Additionally, upper bounds on $\Gamma_p (8n-1)$ have been previously studied by Davis--Mahowald \cite{DM3} (at the prime $2$) and Gonz\'alez \cite{Gonzalez} (at odd primes). At the prime $3$, we will require a novel bound established by the first named author in \Cref{sec:Appendix}.

\begin{thm} \label{thm:gamma-upper}
  We have the following upper bounds on the function $\Gamma_{p} (8n-1)$:
  \begin{enumerate}
    \item (Davis--Mahowald, \cite[Corollary 1.3]{DM3})
      \[\Gamma_2 (8n-1) \leq \frac{3(8n-1)}{10}+7+v_2 (n). \]
    \item (Gonz\'alez, \cite[Theorem 5.1]{Gonzalez}) Assume $p \neq 2$. Then:
      \[\Gamma_p (8n-1) \leq 3+\frac{(2p-1)(8n-1)}{(2p-2)(p^2-p-1)}. \]
    \item (Burklund, \Cref{thm:app-main} (4))\footnote{The statement of \Cref{thm:app-main}(4) in \Cref{sec:Appendix} contains a term depending on a function $\ell(k)$, but this function vanishes by definition when $k=8n-1$.}
      \[\Gamma_3 (8n-1) \leq \frac{25(8n-1)}{184}+19+\frac{1133}{1472}.\]
  \end{enumerate}
\end{thm}
%
%For $p=2$, Davis and Mahowald proved:
%
%\begin{thm} [{\cite[Corollary 1.3]{DM3}}]\label{thm:DM-filt-bound}
%Suppose $\alpha \in \pi_{k} \mathbb{S}_{(2)}$ is in the kernel of the Hurewicz homomorphism 
%$$\pi_{k} \mathbb{S}_{(2)} \to \pi_{k} L_{K(1)} \mathbb{S}.$$
%If the $\mathrm{H}\mathbb{F}_2$-Adams filtration of $\alpha$ is at least
%$$\frac{3k}{10}+4+v_2(k+2)+v_2(k+1),$$
%where $v_2 (a)$ is the $2$-adic valuation of $a$, then $\alpha=0$.
%\end{thm}
%
%\begin{rmk}
%    In the case $k=8n-1$, the $2$-adic valuation terms simplify to \[v_2 ((8n-1)+2) + v_2 ((8n-1)+1) = 3 + v_2 (n).\]
%\end{rmk}
%
%On the other hand, for $p \neq 2$, Jes\'us Gonz\'alez proved
%
%\begin{thm}[{\cite[Theorem 5.1]{Gonzalez}}]\label{thm:G-filt-bound}
%Suppose $\alpha \in \pi_{8n-1} \bS$ is in the kernel of the Hurewicz homomorphism 
%$$\pi_{8n-1} \bS \to \pi_{8n-1} L_{K(1)} \bS.$$
%If the $\mathrm{H}\mathbb{F}_p$ Adams filtration of $\alpha$ is at least
%$$3 + \frac{(2p-1)(8n-1)}{(2p-2)(p^2-p-1)},$$ then $\alpha=0$.
%\end{thm}
%
%Unfortunately, Gonz\'alez's result is insufficient to prove that $w$ is $3$-locally trivial.
%For that reason we have included an appendix which proves the following:
%
%\begin{thm}[Burklund, \Cref{thm:app-main} (4)]\label{thm:B-filt-bound}
%Suppose $\alpha \in \pi_{8n-1} \mathbb{S}_{(3)}$ is in the kernel of the Hurewicz homomorphism 
%$$\pi_{8n-1} \mathbb{S}_{(3)} \to \pi_{8n-1} L_{K(1)} \mathbb{S}.$$
%If the $\mathrm{H}\mathbb{F}_3$-Adams filtration of $\alpha$ is greater than
%%$$\frac{25(8n-1)}{184}+19+\frac{1133}{1472},$$
%then $\alpha=0$.
%\end{thm}
%
%\begin{rmk}
%The statement of \Cref{thm:app-main} in the Appendix contains a term depending on a function $\ell(k)$, but this function vanishes by definition when $k=8n-1$.
%\end{rmk}

\begin{rmk}
\Cref{thm:AdamsBound} and \Cref{thm:app-main} are the only results from the latter half of this paper that are required to settle the Galatius \& Randal--Williams conjecture.  The paper is structured so that the reader willing to assume \Cref{thm:app-main} need not read past \Cref{sec:SyntheticToda} to understand the proof of \Cref{thm:intro-main}.
\end{rmk}

%Since $\pi_{8n-1}\bS_{(p)}$ splits into the direct sum of $(\J_{8n-1})_{(p)}$ and the kernel of the $K(1)$-local Hurewicz homomorphism \cite[VIII.4.2]{EinfBook}, we may safely ignore the difference between the definition of $\Gamma_p (8n-1)$ and the statements involving the Hurewicz map to the $K(1)$-local sphere.
%
At this point, the main work ahead of us in this section is to understand when the bound of \Cref{thm:filt-bound-finale} exceeds the bounds of \Cref{thm:gamma-upper}. To this end, we introduce some compact notation.

%% \begin{ntn}
%%   Let
%%   \begin{align*}
%%     A_2 &:= 2N_2 - 1 = 2h(4n-1) - 2\left\lfloor \log_2(8n) \right\rfloor + 1,
%%     &B_2 := \frac{3(8n-1)}{10} + 7 + v_2(n), \\
%%     A_3 &:= 2N_3 -1 = 2n - 2\left\lfloor \log_3(4n) \right\rfloor - 1
%%     &B_3 := \frac{25(8n-1)}{184} + 19 + \frac{1133}{1472}, \\
%%     A_p &:= 2N_p - 1 = 2\left\lfloor \frac{4n}{2p-2} \right\rfloor - 2\left\lfloor \log_p(4n) \right\rfloor - 1,
%%     &B_p := 3 + \frac{(2p-1)(8n-1)}{(2p-2)(p^2 - p - 1)}.
%%   \end{align*}
%% \end{ntn}

\begin{ntn}
  Let
  $$ A_p := 2N_p - 1\ \ \text{ and }\ \ B_p := \begin{cases} \frac{3(8n-1)}{10} + 7 + v_2(n), & p=2 \\ \frac{25(8n-1)}{184} + 19 + \frac{1133}{1472}, & p=3 \\ 3 + \frac{(2p-1)(8n-1)}{(2p-2)(p^2 - p - 1)}, & p \geq 5 \end{cases}. $$
\end{ntn}

\begin{lem} \label{lem:logs}
  The element $w \in (\pi_{8n-1} \mathbb{S}) / \J_{8n-1}$ is $p$-locally trivial if $A_p > B_p$,
  which occurs for 
  \begin{itemize}
  \item $p=2$ and $n \geq 17$,
  \item $p=3$ and $n \geq 32$,
  \item $p=5$ and $n \geq 16$,
  \item $p=7$ and $n \geq 21$,
  %\item $p=11$ and $n \geq 20$,
  %\item $p=13$ and $n \geq 24$,
  \item $p \ge 11$ and $n \geq 2(2p-2)$.
  \end{itemize}
\end{lem}

\begin{proof}
  The first claim follows from the preceding discussion.
  Verifying the remaining claims is just a matter of checking inequalities between elementary functions.
  For each of these we follow the same basic strategy:
  \begin{enumerate}
  \item Find a smooth function that acts as a lower bound for $A_p - B_p$.
  \item Argue that the derivative of this function is positive.
  \item Find a point where this lower bound is positive.
  \item Go back and fill in small values of $n$ by directly computing $A_p - B_p$.
  \end{enumerate}
  
  % We now prove the first four bullet points. In order to simplify the proof we first introduce some auxiliary functions:
  % \begin{align*}
  %   \overline{A}_p &:= \begin{cases} 4n - 2\log_2 (8n) + 1, & p=2 \\ \frac{4n}{p-1} - 2\log_p (4n) -1, & p\neq 2 \end{cases},
  %                                                           &\overline{B}_p &:= \begin{cases} \frac{3(8n-1)}{10} + 7 + \log_2 (n), & p=2 \\ B_p, & p\neq 2 \end{cases}. 
  % \end{align*}
  We begin with the $p=2$ case.
  \begin{align*}
    A_2 - B_2
    &= 2 \bigg(h(4n-1) - \lfloor \log_2(8n) \rfloor + 1 \bigg) - 1 - \bigg( \frac{3}{10}(8n-1) + 7 + v_2(n) \bigg) \\
    &\geq (4n - 2) - 2\log_2(8n)  + 1 - \frac{3}{10}(8n-1) - 7 - \log_2(n) \\
    &= \frac{8}{5}n - 3\log_2(n)  - 13.7
  \end{align*}
  Since the quantity on the final line is positive for $n=17$ (it is approximately $1.24$) and its derivative with respect to $n$ is positive for $n \geq 3$ we may conclude that $A_2>B_2$ for $n \geq 17$.
  % The inequality for the remaining value of $n$ can now be read off from \Cref{tbl:big}.

  Next we handle the $p=3$ case, which proceeds in the same fashion as $p=2$.
  \begin{align*}
    A_3 - B_3
    &= 2 \left( \left\lfloor \frac{4n}{4} \right\rfloor - \lfloor \log_3(4n) \rfloor \right) - 1 - \left( \frac{25}{184}(8n-1) + 19 + \frac{1133}{1472} \right) \\
    &\geq 2n - 2\log_3(4n) - 1 - \frac{200}{184}n - 20 \\
    &= \frac{21}{23}n - 2\log_3(4n) - 21      
  \end{align*}
  The quantity on the final line is positive for $n=33$ (it is approximately $0.24$) and its derivative with respect to $n$ is positive for $n \geq 2$, thus we may conclude that $A_3>B_3$ for $n \geq 33$.
  The inequality for the remaining value of $n$ can now be read off from \Cref{tbl:big}.
  % Again, the remaining cases can be read off from \Cref{tbl:big}.

  % \begin{table}
  %   \begin{center}
  %     % \renewcommand{\arraystretch}{1.2}
  %     \begin{tabular}{|c|c||c|c|c||c|c||c|c||c|c||c|c|c|}\hline
  %       $n$ & $8n-1$ & $A_2$ & $B_2$ & $\overline{B}_2$ & $A_3$ & $B_3$ & $A_5$ & $B_5$ & $A_7$ & $B_7$ \\\hline\hline
  %         16 & 127 & 49 & 49.1 & 49.10 & 25 & 37.03 & 11 & 10.52 & 5 & 6.36 \\\hline
  %         17 & 135 & 55 & 47.50 & 51.59 & 27 & 38.11 & 11 & 10.99 & 5 & 6.57 \\\hline
  %         18 & 143 & 57 & 50.90 & 54.07 & 29 & 39.20 & 13 & 11.47 & 7 & 6.78 \\\hline
  %         19 & 151 & 63 & 52.30 & 56.55 & 31 & 40.29 & 13 & 11.94 & 7 & 6.99 \\\hline
  %         20 & 159 & 65 & 56.70 & 59.02 & 33 & 41.37 & 15 & 12.41 & 7 & 7.20 \\\hline
  %         21 & 167 & 71 & 57.10 & 61.49 & 33 & 42.46 & 15 & 12.89 & 9 & 7.41 \\\hline
  %         22 & 175 & 73 & 60.50 & 63.96 & 35 & 43.55 & 17 & 13.36 & 9 & 7.62 \\\hline
  %         23 & 183 & 79 & 61.90 & 66.42 & 37 & 44.63 & 17 & 13.84 & 9 & 7.84 \\\hline
  %         24 & 191 & 81 & 67.30 & 68.88 & 39 & 45.72 & 19 & 14.31 & 11 & 8.05 \\\hline
  %         25 & 199 & 87 & 66.70 & 71.34 & 41 & 46.81 & 19 & 14.78 & 11 & 8.26 \\\hline
  %         26 & 207 & 89 & 70.10 & 73.80 & 43 & 47.89 & 21 & 15.26 & 11 & 8.47 \\\hline
  %         27 & 215 & 95 & 71.50 & 76.25 & 45 & 48.98 & 21 & 15.73 & 13 & 8.68 \\\hline
  %         28 & 223 & 97 & 75.90 & 78.71 & 47 & 50.07 & 23 & 16.20 & 13 & 8.89 \\\hline
  %         29 & 231 & 103 & 76.30 & 81.16 & 49 & 51.16 & 23 & 16.68 & 13 & 9.10 \\\hline
  %         30 & 239 & 105 & 79.70 & 83.61 & 51 & 52.24 & 25 & 17.15 & 15 & 9.32 \\\hline
  %         31 & 247 & 111 & 81.10 & 86.05 & 53 & 53.33 & 25 & 17.62 & 15 & 9.53 \\\hline
  %         32 & 255 & 111 & 88.50 & 88.50 & 55 & 54.42 & 25 & 18.10 & 15 & 9.74 \\\hline
  %         33 & 263 & 117 & 85.90 & 90.94 & 57 & 55.50 & 25 & 18.57 & 17 & 9.95 \\\hline
  %         34 & 271 & 119 & 89.30 & 93.39 & 59 & 56.59 & 27 & 19.05 & 17 & 10.16 \\\hline
  %         35 & 279 & 125 & 90.70 & 95.83 & 61 & 57.68 & 27 & 19.52 & 17 & 10.37 \\\hline
  %         36 & 287 & 127 & 95.10 & 98.27 & 63 & 58.76 & 29 & 19.99 & 19 & 10.58 \\\hline
  %     \end{tabular}
  %     \renewcommand{\arraystretch}{1.0}
  %       \vspace{0.5cm}
  %       \caption{}
  %       \label{tbl:big}
  %   \end{center}
  % \end{table}

  %% old version of table 
  % \begin{table}
  %   \begin{center}
  %     \renewcommand{\arraystretch}{1.1}
  %     \begin{tabular}{|c||c|c||c|c||c|c|}\hline
  %       $n$ & $A_2$ & $B_2$ & $A_5$ & $B_5$ & $A_7$ & $B_7$ \\\hline\hline
  %         16 & 49 & 49.1 &  11 & 10.52 & 5 & 6.36 \\\hline
  %         17 & 55 & 47.50 &  11 & 10.99 & 5 & 6.57 \\\hline
  %         18 & 57 & 50.90 &  13 & 11.47 & 7 & 6.78 \\\hline
  %         19 & 63 & 52.30 &  13 & 11.94 & 7 & 6.99 \\\hline
  %         20 & 65 & 56.70 &  15 & 12.41 & 7 & 7.20 \\\hline
  %         21 & 71 & 57.10 &  15 & 12.89 & 9 & 7.41 \\\hline
  %         22 & 73 & 60.50 &  17 & 13.36 & 9 & 7.62 \\\hline
  %         23 & 79 & 61.90 &  17 & 13.84 & 9 & 7.84 \\\hline
  %         24 & 81 & 67.30 &  19 & 14.31 & 11 & 8.05 \\\hline
  %     \end{tabular}
  %     \hspace{0.5cm}
  %     \begin{tabular}{|c||c|c|}\hline
  %       $n$ & $A_3$ & $B_3$ \\\hline\hline
  %       30 & 51 & 52.24 \\\hline
  %       31 & 53 & 53.33 \\\hline
  %       32 & 55 & 54.42 \\\hline
  %       33 & 57 & 55.50 \\\hline
  %       %34 & 59 & 56.59 \\\hline
  %       %35 & 61 & 57.68 \\\hline
  %       %36 & 63 & 58.76 \\\hline
  %     \end{tabular}
  %     \renewcommand{\arraystretch}{1.0}
  %       \vspace{0.5cm}
  %       \caption{}
  %       \label{tbl:big}
  %   \end{center}
  % \end{table}

  %% p=2 removed from table
  \begin{table}
    \begin{center}
      \renewcommand{\arraystretch}{1.1}
      \begin{tabular}{|c||c|c|}\hline
        $n$ & $A_3$ & $B_3$ \\\hline\hline
        % 30 & 51 & 52.24 \\\hline
        31 & 53 & 53.33 \\\hline
        32 & 55 & 54.42 \\\hline
        33 & 57 & 55.50 \\\hline
        %34 & 59 & 56.59 \\\hline
        %35 & 61 & 57.68 \\\hline
        %36 & 63 & 58.76 \\\hline
      \end{tabular}
      \hspace{0.5cm}
      \begin{tabular}{|c||c|c|}\hline
        $n$ & $A_5$ & $B_5$  \\\hline\hline
          16 &   11 & 10.52  \\\hline
          17 &   11 & 10.99  \\\hline
          18 &   13 & 11.47  \\\hline
          19 &   13 & 11.94  \\\hline
          20 &   15 & 12.41  \\\hline
          21 &   15 & 12.89  \\\hline
          22 &   17 & 13.36  \\\hline
          23 &   17 & 13.84  \\\hline
          24 &   19 & 14.31  \\\hline
      \end{tabular}
      \hspace{0.5cm}
      \begin{tabular}{|c||c|c|}\hline
        $n$ & $A_7$ & $B_7$ \\\hline\hline
        20 &  7 & 7.20 \\\hline
        21 &  9 & 7.41 \\\hline
        22 &  9 & 7.62 \\\hline
        23 &  9 & 7.84 \\\hline
        24 &  11 & 8.05 \\\hline
      \end{tabular}
      \hspace{0.5cm}
      \renewcommand{\arraystretch}{1.0}
        \vspace{0.5cm}
        \caption{}
        \label{tbl:big}
    \end{center}
  \end{table}

  We handle the remaining values of $p$ uniformly.
  Let $(p-1)k = n$, then
  \begin{align*}
    A_p - B_p &= 2 \left( \left\lfloor \frac{4n}{2p-2} \right\rfloor - \lfloor \log_{p}(4n) \rfloor \right) - 1 - \left( 3 + \frac{(2p-1)(8n-1)}{(2p-2)(p^2-p-1)} \right) \\
    &\geq 2 \frac{4n}{2p-2} - 2\log_{p}(4n) - 6 - \frac{(2p-1)8n}{(2p-2)(p^2-p-1)} \\
    &\geq 4k - 2\log_{p}(4(p-1)k) - 6 - \frac{(2p-1)4k}{p^2 - p - 1} \\
    &\geq 4k - \frac{(2p-1)4k}{p^2 - p - 1} - 8 - 2\log_{p}(4k). 
  \end{align*}
  Let $C_p$ denote the quantity on the final line.
  % As above we will conclude that $C_p$ is positive using positivity of derivatives together with positivity at specific points.
  For $p \geq 5$ and $k \geq 2$ we have that
  \begin{align*}
    \frac{\partial}{\partial k} C_p
    &= 4 - \frac{(2p-1)4}{p^2 - p - 1} - \frac{2}{\log(p) k} \geq 4 - \frac{8p}{p^2 - 2p} - \frac{2}{\log(p) k} \\
    &\geq 4 - \frac{8}{3} - \frac{2}{2\log(5)} > 0
  \end{align*}  
  and
  \[ \frac{\partial}{\partial p} C_p = 4k\frac{2p^2 - 2p + 3}{(p^2 - p - 1)^2} + \frac{2\log(4k)}{(\log(p))^2 p} = \frac{2k((2p - 1)^2 + 5)}{(p^2 - p - 1)^2} + \frac{2\log(4k)}{(\log(p))^2 p}> 0. \]
  When $p=5$ and $k=6$ the value of $C_p$ is approximately $0.68$ and
  when $p=7$ and $k=4$ the value of $C_p$ is approximately $0.08$.
  Translating back into terms of $p$ and $n$ instead of $k$ this completes the proof of the lemma for $p \geq 11$ and all but finitely many cases for $p=5,7$. 
  A second consultation with \Cref{tbl:big} now completes the proof.
\end{proof}

\begin{proof}[Proof of Theorem \ref{thm:w=0}]
  
  In order to finish the proof of \Cref{thm:w=0} at $p\geq 5$ it will suffice to show that $\pi_{8n-1}\Ss^0_{(p)}$ is generated by the image of $J$ for each $n$ below the bound from \Cref{lem:logs}.
  Through this range the $\mathrm{E}_2$-page of the Adams--Novikov spectral sequence is calculated in \cite[Theorem 4.4.20]{GreenBook} and the spectral sequence degenerates at the $\mathrm{E}_2$-page for degree reasons.

  In \Cref{tbl:small-stems} we list the generators of the cokernel of $J$ in degrees below the bound from \Cref{lem:logs} at primes $5,7,11$ and $13$.
  None of these generators are in a degree congruent to $-1$ modulo 8.
  At primes $p \geq 17$ we argue as follows:
  again using \cite[Theorem 4.4.20]{GreenBook} we know that the first element of $\coker(J)$ at odd primes is $\beta_1$ in the $(2p^2 - 2p -2)$ stem.
  On the other hand, since
  \[8n-1 \leq 16(2p-2) - 1 < 2p^2 - 2p - 2 \]
  the bound from \Cref{lem:logs} is below this point.
\end{proof}
  
  % To do this we make use of the known structure of the homotopy groups of the $p$-local sphere in small (relative to $p$) degree, which is summarized in \cite[Theorem 4.4.20]{GreenBook}.

  \begin{table}
    \begin{center}
      \renewcommand{\arraystretch}{1.1}
      \begin{tabular}{|c|c|c|c|}\hline
        $p$ & Max $n$ & Max degree & Classes in $\coker(J)$ \\\hline\hline
        5 & 15 & 119 & $ \beta_1,\ \alpha_1\beta_1,\ \beta_1^2,\ \alpha_1\beta_1^2,\ \beta_2,\ \alpha_1\beta_2,\ \beta_1^3 $\hfill\vadjust{} \\\hline
        7 & 20 & 159 & $\beta_1,\ \alpha_1\beta_1$\hfill\vadjust{} \\\hline
        11 & 39 & 311 & $\beta_1,\ \alpha_1\beta_1$\hfill\vadjust{} \\\hline
        13 & 47 & 375 & $\beta_1,\ \alpha_1\beta_1$\hfill\vadjust{} \\\hline             
      \end{tabular}
      \renewcommand{\arraystretch}{1.0}
      \vspace{0.5cm}
      \caption{}
      \label{tbl:small-stems}
    \end{center}
  \end{table}
  
%   Suppose that $p = 5$.
%   The maximum value of $8n-1$ allowed is 119.
%   Calculation of the low-dimensional homotopy of the $5$-local sphere informs us that through this range $\coker(J)$ has 7 generators,
% $$ \beta_1,\ \alpha_1\beta_1,\ \beta_1^2,\ \alpha_1\beta_1^2,\ \beta_2,\ \alpha_1\beta_2, \text{ and } \beta_1^3. $$
% None of these are in a degree congruent to $-1$ modulo 8.
% The same argument goes through for $p = 7,11,$ and $13$.

%% Suppose that $p = 7$.
%% The maximum value of $8n-1$ allowed is 159.
%% Through this range $\coker(J)$ has two generators $\beta_1$ and $\alpha_1\beta_1$.
%% Neither is in a degree congruent to $-1$ modulo 8.

%% Suppose that $p = 11$.
%% The maximum value of $8n-1$ allowed is 151.
%% Through this range $\coker(J)$ is empty.

%% Suppose that $p = 13$.
%% The maximum value of $8n-1$ allowed is 183.
%% Through this range $\coker(J)$ is empty.


% It is clear that, for $p \ge 5$, the above corollary will apply for all sufficiently large $n$.  On the other hand, as summarized for example in \cite[4.4.20]{GreenBook}, the low-dimensional homotopy groups of the $p$-local sphere spectrum are well-understood, and we can make direct arguments for small values of $n$.  In particular, combining \cite[4.4.22]{GreenBook} and \cite[4.4.20]{GreenBook} gives a complete description of $\pi_*(\bS_{(p)})$ for $* \le (p^2+p)(2p-2)$. \todo{add a few more details}

\subsection{A note on the remaining dimensions}
\label{sec:open-problems}\ 

Galatius and Randal--Williams conjecture \cite[Conjecture A]{GRAbelianQuotients} that the map
$$\pi_{8n-1} \bS \to \pi_{8n-1} \MO \langle 4n \rangle$$
has kernel equal to $\mathcal{J}_{8n-1}$ for all $n \ge 1$, and not just for $n >31$.

It is known that the conjecture is true when $n=1$ and when $n=2$ \cite[p. 13]{GRAbelianQuotients}.  Indeed, the case $n=1$ follows from the fact that there is nothing in $\pi_7 \bS$ not in the image of $J$.  When $n=2$, it follows from direct calculation of $\pi_{15} \MO \langle 8 \rangle=\pi_{15} \mathrm{MString}$ as in \cite{GiambalvoO8,GiambalvoO8corrected}.  The methods of this paper first apply when $n \ge 3$.

The first and third named authors returned to this question in \cite{ManSmall} and proved the following theorem, completely resolving the remaining cases of the Galatius and Randal--Williams conjecture:
\begin{thm}[{\cite[Theorem 1.4]{ManSmall}}]
  The kernel of the map
  \[\pi_{2k-1} \bS \to \pi_{2k-1} \MO \langle k \rangle\]
  is equal to $\mathcal{J}_{2k-1}$ if $k \neq 9,12$.
  When $k=9$, it is generated by $\mathcal{J}_{17}$ and $\eta \eta_4 \in \pi_{17} (\bS)$.
  When $k=12$, it is generated by $\mathcal{J}_{23}$ and $\eta^3 \overline{\kappa} \in \pi_{23} (\bS)$.

  In particular, the Galatius and Randal--Williams conjecture is true precisely when $n \neq 3$.
\end{thm}
%
%It seems worthwhile to comment on the prospect of settling the conjecture in general, since this would expand the range of dimensions in which we could conclude the applications of Sections \ref{sec:classification} and \ref{sec:Applications}.
%
%It is known that the conjecture is true when $n=1$ and when $n=2$ \cite[p. 13]{GRAbelianQuotients}.  Indeed, the case $n=1$ follows from the fact that there is nothing in $\pi_7 \bS$ not in the image of $J$.  When $n=2$, it follows from direct calculation of $\pi_{15} \MO \langle 8 \rangle=\pi_{15} \mathrm{MString}$ as in \cite{GiambalvoO8,GiambalvoO8corrected}.  The methods of this paper first apply when $n \ge 3$.
%
%For many values of $n$, our ability to prove the Galatius and Randal--Williams conjecture is limited by our knowledge of the $2$-primary and $3$-primary Adams spectral sequences in `low dimensions.'  We suspect that, as knowledge of the low-dimensional homotopy groups of spheres is pushed forward, the following conjecture will eventually be verified:
%
%\begin{cnj} \label{cnj:GRWConj}
%The Galatius and Randal--Williams conjecture is true for $n \ne 3,4,6$.  More specifically, given knowledge of the $2$-primary and $3$-primary Adams spectral sequences for the sphere, the conjecture for $n \ne 3,4,6$ follows from our Theorem \ref{thm:filt-bound-finale}.
%\end{cnj}
%
%For the convenience of the reader interested in this conjecture we provide \Cref{tbl:remaining}, which tabulates our bounds on the Adams filtration of $w$ in all remaining cases. Even given perfect knowledge of Adams spectral sequences, which we have in dimensions less than $60$ \cite{Isaksen}, our methods fail to settle the Galatius and Randal--Williams conjecture for $n=3,4$ or $6$.  Our methods come close to settling the conjecture for $n=4$ and for $n=6$.  On the other hand, Theorem \ref{thm:filt-bound-finale} has little content when $n=3$.  In fact, preliminary computations of the first author and Zhouli Xu indicate that the Galatius and Randal--Williams conjecture is false for $n=3$.  Specifically, computations suggest that the kernel of
%$$\pi_{23} \mathbb{S} \to \pi_{23} \MO \langle 12 \rangle$$
%contains $\eta^3 \bar{\kappa}$, which is not in $\mathcal{J}_{23}$.

%% \begin{comment}
%% \begin{cnj}
%% The Galatius and Randal--Williams conjecture is false for $n=3$.  In other words, the map
%% $$\pi_{23} \mathbb{S} \to \pi_{23} \MO \langle 12 \rangle$$
%% has kernel larger than $\mathcal{J}_{23}$.
%% \end{cnj}
%% \end{comment}


%\begin{table}
%    \begin{center}
      % \renewcommand{\arraystretch}{1.2}
%      \scalebox{0.95}{
%      \begin{tabular}{|c||c|c|c|c|c|c|c|c|c|c|c|c|c|c|c|} \hline
%        $n$  & 3 & 4 & 5 & 6 & 7 & 8 & 9 & 10 & 11 & 12 & 13 & 14 & 15 & 16 & 17  \\\hline
%        $8n-1$  & 23 & 31 & 39 & 47 & 55 & 63 & 71 & 79 & 87 & 95 & 103 & 111 & 119 & 127 & 135 \\\hline
%        $2N_2 - 1$  & 5 & 5 & 11 & 13 & 19 & 19 & 25 & 27 & 33 & 35 & 41 & 43 & 49 & 49 & 55 \\\hline
%        $2N_3 - 1$  & 1 & 3 & 5 & 7 & 7 & 9 & 11 & 13 & 15 & 17 & 19 & 21 & 23 & 25 & 27 \\\hline
%      \end{tabular}}
%      % \renewcommand{\arraystretch}{1.0}
%      \vspace{0.3cm}
%
%      \scalebox{0.95}{
%      \begin{tabular}{|c||c|c|c|c|c|c|c|c|c|c|c|c|c|c|} \hline
%        $n$  & 18 & 19 & 20 & 21 & 22 & 23 & 24 & 25 & 26 & 27 & 28 & 29 & 30 & 31 \\\hline
%        $8n-1$  & 143 & 151 & 159 & 167 & 175 & 183 & 191 & 199 & 207 & 215 & 223 & 231 & 239 & 247 \\\hline
%        $2N_3 - 1$  & 29 & 31 & 33 & 33 & 35 & 37 & 39 & 41 & 43 & 45 & 47 & 49 & 51 & 53 \\\hline
%      \end{tabular}}
%      \vspace{0.3cm}
%      \caption{}
%      \label{tbl:remaining}
%    \end{center}
%\end{table}


%%% Local Variables:
%%% mode: latex
%%% TeX-master: "Main"
%%% eval: (visual-line-mode t)
%%% eval: (auto-fill-mode 0)
%%% End:
